\documentclass[11pt]{article}
\usepackage[T1]{fontenc}
\usepackage{lmodern,microtype,amsmath,amsthm,amssymb,mathtools,booktabs}
\usepackage[letterpaper,left=1.01in,right=1.01in,top=.87in,bottom=.86in]{geometry}
\usepackage{xcolor,hyperref}
\definecolor{ink}{HTML}{194F67}
\hypersetup{colorlinks=true,urlcolor=ink,linkcolor=ink,citecolor=ink,pdftitle={Squarefree products of small primes after arbitrary deletions}}
\newtheorem{theorem}{Theorem}[section]
\newtheorem{lemma}[theorem]{Lemma}
\newtheorem{proposition}[theorem]{Proposition}
\newtheorem{corollary}[theorem]{Corollary}
\newtheorem{remark}[theorem]{Remark}
\DeclareMathOperator{\Prob}{\mathbb P}
\renewcommand{\epsilon}{\varepsilon}
\setlength{\emergencystretch}{2em}
\title{\textbf{Squarefree products of small primes\protect\\ after arbitrary deletions}}
\author{Research note}
\date{26 November 2026}
\begin{document}
\maketitle
\begin{abstract}
Every unit modulo a sufficiently large integer $q$ is a product of exactly
$\lceil B\log q\rceil$ distinct primes below $C(\log q)^2$.
This still holds, with uniform asymptotic representation counts, after deletion
of any fraction at most $1/2-\delta$ of the eligible primes. For any fixed
retained density $a>0$ we identify a quotient of order at most $\lfloor1/a\rfloor$
whose exact subset-product distribution determines both support and relative
counts in the full group. This gives the smoothness scale proposed by Booker
and Pomerance for squarefree representatives. The analytic input for every
modulus is GRH for Dirichlet characters.
For primes up to $(\log q)^A$, $A>2$, a moment argument reduces a sufficient
distance from either endpoint of the subset-size range to order
$\log q/\log\log q$. A large-sieve argument gives an unconditional
density-one version for prime moduli. A common list of $O(\log Q)$ primes
suffices for every modulus in $[Q,2Q]$ under Dirichlet GRH.
\end{abstract}

\section{Squarefree representatives at the logarithm-squared scale}
A positive integer is $y$-smooth if all its prime factors are at most $y$.
Booker and Pomerance asked how small $y$ can be if every nonzero class modulo
a prime $p$ must contain a squarefree $y$-smooth integer. In their final
section they suggested that GRH for Kummerian fields should give
$y=O((\log p)^2)$ \cite{bp}. Squarefreeness matters here. A set of primes
can generate the unit group by ordinary multiplication while exhausting its
available elements before a squarefree product reaches a desired class.

Put $G=(\mathbb Z/q\mathbb Z)^\times$, $L=\log q$,
$P=\{p\le y:p\text{ prime},\ p\nmid q\}$, and $N=|P|$. All logarithms
are natural. For $S\subseteq P$ put $m=|S|$. Choose a $k$-element subset of
$S$ uniformly and let $X_k$ denote its product in $G$.
Thus $\binom mk\Prob(X_k=g)$ counts distinct positive squarefree integers,
each with exactly $k$ prime factors. All estimates below are uniform in the
chosen primes, target residue, and indicated subset sizes.

\begin{theorem}[Majority after deletions]\label{thm:major}
Assume Dirichlet GRH. Given $0<\delta<1/2$ there are constants $B,C>0$ such that,
for all sufficiently large $q$, with $y=CL^2$ and $m\ge(1/2+\delta)N$,
\[
 \#\{T\subseteq S: |T|=k,\ \textstyle\prod_{p\in T}p\equiv g\pmod q\}
 =\frac{\binom mk}{\varphi(q)}(1+O(q^{-2}))
 \qquad(g\in G)
\]
for every integer $BL\le k\le m-BL$.
\end{theorem}
In particular one may take $S=P$ and $k=\lceil B\log q\rceil$, with
absolute $B,C$. This proves the proposed upper bound for the least smoothness
parameter using the Dirichlet case of GRH. The theorem also prescribes the
number of factors and allows their pool to be chosen adversarially. For the
finitely many small prime moduli the smoothness constant can be enlarged,
since Dirichlet's theorem supplies a prime in each nonzero class.

\section{The quotient that determines all obstructions}
For $\chi\in\widehat G$ and nonempty $U\subseteq P$, write
\[
 b_\chi(U)=\left|\frac1{|U|}\sum_{p\in U}\chi(p)\right|.
\]
For $0<\tau<1$ define the subgroup of the character group
\[
 H=\langle\chi\in\widehat G:b_\chi(S)>1-\tau\rangle,
 \qquad K=\bigcap_{\chi\in H}\ker\chi,\qquad Q=G/K .
\]
The identity character belongs to $H$. Finite abelian duality gives
$\widehat Q=H$ and $d:=|Q|=|H|$. Denote the quotient map by $\pi$ and the
probability that $\pi(X_k)=v$ by $\lambda_k(v)$.

\begin{theorem}[Exact support and relative counts]\label{thm:main}
Fix $0<a\le1$ and put $D=\lfloor1/a\rfloor$. There are
$\tau>0$ and $B,C>0$ depending only on $a$ for which the following holds under
Dirichlet GRH. Take $y=C L^2$, $m\ge aN$ and $BL\le h:=\min(k,m-k)$.
For all sufficiently large $q$, the same subgroup $H$ works for all these $k$,
$d\le D$, and
\begin{equation}\label{eq:profile}
 \left|\varphi(q)\Prob(X_k=g)-d\lambda_k(\pi(g))\right|\le q^{-2}
 \quad (g\in G).
\end{equation}
Every fiber of $\pi$ that contains one product contains every residue as a
product. On every such fiber,
\[
 \Prob(X_k=g)=\frac{d\lambda_k(\pi(g))}{\varphi(q)}(1+o(1)).
\]
The support is exactly $\pi^{-1}(\operatorname{supp}\lambda_k)$.
\end{theorem}
The quotient has bounded order even though $m$, $k$, $q$, and $\varphi(q)$ all
grow. Its exact probabilities can be computed by multiplying
$\prod_{p\in S}(1+z[\pi(p)])$ in the group ring $\mathbb Z[Q][z]$ and
dividing the degree-$k$ counts by $\binom mk$. Here $[u][v]=[uv]$.

\paragraph{The half and third thresholds.}
For $a>1/2$, $d=1$, proving Theorem~\ref{thm:major} from
Theorem~\ref{thm:main}. For $a>1/3$, the only other possibility is $d=2$.
In this case one nonprincipal quadratic character gives the quotient. If $S$
uses both signs then $1\le k<m$ permits two $k$-sets differing by a swap of
opposite signs. Both quotient values occur, so every unit is represented.
If $S$ lies in a single quadratic sign class, its products lie in the one sign
class fixed by $k$. Equivalently, for $a>1/3$, full coverage in the theorem's
range holds exactly when no nonprincipal quadratic character is constant on
$S$. If $s$ of the $m$ primes have negative sign, the two quotient masses are
$(1\pm E)/2$ with
\[
 E=\binom mk^{-1}[z^k](1+z)^{m-s}(1-z)^s.
\]
Even one prime of the opposite sign suffices for coverage, although the
counts of the two signs can differ markedly.

\paragraph{Sharp quotient size.}
Given an integer $d\ge2$, take primes $q\equiv1\pmod d$ and a quotient of
$G$ of order $d$. Retain the most populous coset among the primes in $P$.
This retains at least $N/d$ primes and confines all products of fixed length
to one coset. Any quotient that exactly explains this nonempty support must
have order at least $d$, since its kernel must lie in the kernel of the chosen
quotient. Thus reaching density $1/d$ genuinely permits an obstruction of
order $d$. In particular one cannot lower the majority guarantee to one
half. At density $1/3$ a cubic coset supplies an obstruction of order three.

\section{Two character estimates for the full prime pool}
The proof requires an estimate uniform over conductors dividing $q$. The
ordinary sharp-cutoff bound under Dirichlet GRH is
\begin{equation}\label{eq:sharp}
 \left|\sum_{p\le x}\chi(p)\right|\ll \sqrt x\log(qx)
 \qquad(\chi\ne1).
\end{equation}
Here and below a character vanishes on nonunits. Equation~\eqref{eq:sharp}
is Montgomery--Vaughan, Theorem 13.7, equation (13.21) \cite{mv}. It yields
a power saving in $\log q$ for $x=L^A$, $A>2$. At $x=CL^2$ we use their
once-integrated estimate, equation (13.28),
\begin{equation}\label{eq:tri}
 F_\chi(t):=\sum_{n\le t}(t-n)\Lambda(n)\chi(n)
 \ll t^{3/2}\log(2q),\qquad t\ge2,\quad\chi\ne1 .
\end{equation}
Its inverse Mellin kernel includes the factor $1/(s(s+1))$. Each critical-line
zero contributes $t^{\rho+1}/(\rho(\rho+1))$, with an absolutely summable
weight in its ordinate. This is where the square-root estimate survives
without a factor $\log t$. Both cited character bounds apply to imprimitive
characters as well.

\begin{lemma}[Small bias at the endpoint]\label{lem:analytic}
Given $\epsilon>0$, one can choose a constant $C$ so that, for $y=C L^2$
and all sufficiently large $q$,
$\max_{\chi\ne1} b_\chi(P)\le\epsilon$.
In both regimes $y=C L^2$ and $y=L^A$, $A>2$, we have $N\sim y/\log y$.
\end{lemma}
\begin{proof}
The estimate for $N$ follows from the prime number theorem after deleting at
most $L/\log 2$ primes dividing $q$. Fix $0<\eta<1/2$. The weight
\[
 w(n)=\frac{(y-n)_+-((1-\eta)y-n)_+}{\eta y}
\]
equals one up to $(1-\eta)y$ and tapers linearly to zero at $y$.
Equation~\eqref{eq:tri} bounds $|\sum_n w(n)\Lambda(n)\chi(n)|$ by
$O(\sqrt y\log(2q)/\eta)$. Removing prime powers of exponent at least two
costs $O(\sqrt y\log^2 y)=o(y)$. Restoring the last $\eta y$ of the interval
costs at most $\sum_{(1-\eta)y<p\le y}\log p=(\eta+o(1))y$ by the prime
number theorem with fixed $\eta$. Finally,
\[
 \sum_{p\le y}(\log y-\log p)=\pi(y)\log y-\vartheta(y)=o(y).
\]
Thus replacing $\log p$ by $\log y$ costs $o(y)$ even after character
twisting. Division by $N\log y\sim y$ gives
$b_\chi(P)\le \eta+O(1/(\eta\sqrt C))+o(1)$. Choose first $\eta$ and then
$C$. All errors are uniform in $\chi$.
\end{proof}

\section{Extraction of a quotient of optimal order}
Work first with a finite pool $P$ of labeled group elements with
$b_\chi(P)\le\epsilon$ for all nonprincipal $\chi$. We use the probability
inner product on functions on $P$. Distinct character vectors have inner
product of absolute value at most $\epsilon$. For any $t$ such vectors
their Gram matrix has largest eigenvalue at most $1+(t-1)\epsilon$ by
the row-sum bound. Apply the resulting Bessel inequality to the indicator
of $S$ and put $\alpha=m/N$. If $b_\chi(S)\ge u$ for each of these $t$
characters, then
\[
 t\alpha^2 u^2\le\alpha(1+(t-1)\epsilon),\qquad
 t\le\frac{1-\epsilon}{\alpha u^2-\epsilon}
 \quad (\alpha u^2>\epsilon).
\tag{G}\label{eq:gram}
\]
On $S$ a character of bias $>1-\tau$ is within $\sqrt{2\tau}$, in
normalized $L^2$, of a unit constant. Its inverse has the same property.
By telescoping a product and using the triangle inequality, any word of
length at most $D$ in these generators and inverses is within
$D\sqrt{2\tau}$ of some unit constant. Its bias is at least
$1-D^2\tau$ by expanding squared distance. Since $1/a<D+1$, we may choose
$\tau>0$, $1-D^2\tau>0$, and then $\epsilon>0$ so small that
\[
 \frac{1-\epsilon}{a(1-D^2\tau)^2-\epsilon}<D+1
\]
with positive denominator. The radius-$D$ word ball in $H$ contains at most
$D$ elements by \eqref{eq:gram}. Starting at the identity, a word ball
gains at least one new element at each step until it is the whole subgroup.
Consequently $|H|\le D$. By construction,
$b_\chi(S)\le1-\tau$ for every $\chi\notin H$.

\section{Fixed-cardinality Fourier inversion}
For $z_p=\chi(p)$ put
\[
 E_j(\chi)=\binom mj^{-1}[z^j]\prod_{p\in S}(1+zz_p)
          =\mathbb E[\chi(X_j)] .
\]
Taking complements of subsets gives
$E_{m-j}(\chi)=(\prod_{p\in S}z_p)\overline{E_j(\chi)}$.
For $1\le h\le m/2$, put $r=h/(m-h)$ and $b=b_\chi(S)$.
For $|w|=1$,
\[
 \frac{|1+rw|}{1+r}
 \le\exp\!\left(-\frac{r(1-\Re w)}{(1+r)^2}\right),
\]
by squaring both sides and using $1-u\le e^{-u}$. Multiplication on the
circle of radius $r$ and Cauchy's coefficient estimate give
\begin{equation}\label{eq:circle}
 |E_k(\chi)|=|E_h(\chi)|\le (m+1)
       \exp\{-h(1-h/m)(1-b)\},\qquad h=\min(k,m-k).
\end{equation}
Indeed the normalization costs the reciprocal of the probability of $h$ in
$\operatorname{Bin}(m,h/m)$, which is at most $m+1$ since $h$ is a mode.
Fourier inversion in $G$ and in $Q$ now gives the exact bound
\begin{equation}\label{eq:inv}
 |\,|G|\Prob(X_k=g)-d\lambda_k(\pi(g))\,|
 \le\sum_{\chi\notin H}|E_k(\chi)|.
\end{equation}
The characters that factor through $Q$ give exactly $d\lambda_k(\pi(g))$.
At the endpoint, the remaining sum is at most
$q(m+1)\exp(-\tau h/2)$. Since $m\le CL^2$, taking, for example,
$B=10/\tau$ gives \eqref{eq:profile}.

\begin{lemma}[A reachable quotient value has appreciable mass]\label{lem:swap}
For any finite abelian quotient of order $d$, every nonzero value of the
probability mass function $\lambda_k$ is at least
\[
 \left(\sum_{i=0}^{d-1}\binom ki\binom{m-k}i\right)^{-1}
 \ge (d\,m^{2(d-1)})^{-1}
\]
for $m\ge1$. Out-of-range binomial coefficients are zero.
\end{lemma}
\begin{proof}
Fix a target subset $T$ with the desired quotient product. For any other
$k$-subset $U$, pair $U\setminus T$ bijectively with $T\setminus U$.
Each pair is a possible swap with a specified quotient ratio. Their joint
product corrects the quotient of $U$ to the target. In a list of $d$ or more
ratios two of the partial products, including the empty product, coincide.
Remove that nonempty identity-product block and repeat. At most $d-1$ swaps
remain with the same total correction. Every $U$ is therefore within $d-1$
swaps of some target subset. The sum in the denominator is the size of
one such neighborhood. Counting their union proves the result.
\end{proof}
If the quotient mass is positive, it is bounded below by an inverse power
of $\log q$ since $d$ is fixed and $m$ is polylogarithmic.
The error in \eqref{eq:profile} is smaller, which yields the relative estimate
and the exact support in Theorem~\ref{thm:main}. A zero quotient mass
excludes the entire fiber by definition.
\section{Polynomial-size representatives from distinct primes}
\label{sec:acceleration}

Suppose that $A>2$, $L=\log q$, and $P$ consists of the primes at most
$y=L^A$ that are coprime to $q$. Write $N=|P|$, take $S\subseteq P$
with $m=|S|\ge aN$ for fixed $a>0$, and put
\[
 c_A=\frac{2(A-1)}{A(A-2)}.
\]
Use the subgroup $H$ of the character group constructed above, generated
by characters with $b_\chi(S)>1-\tau$, where
$b_\chi(U)=|\sum_{p\in U}\chi(p)|/|U|$. Denote the normalized $k$th
elementary symmetric coefficient of $(\chi(p))_{p\in S}$ by $E_k(\chi)$.

\begin{theorem}[Short products]
\label{thm:short-products}
Assume Dirichlet GRH. For every fixed $c>c_A$, there is a constant
$\gamma>0$, independent of $q,S,k$, such that, for all sufficiently large $q$,
\[
 h=\min(k,m-k)\ge c\frac{L}{\log L}
 \quad\Longrightarrow\quad
 \sum_{\chi\notin H}|E_k(\chi)|\le
 \exp\!\left(-\gamma\frac{L}{\log L}\right).
\]
Consequently the same quotient law and exact support description hold
on this shorter scale. With $k=\lceil cL/\log L\rceil$ every attainable
class has a positive squarefree representative at most $q^{Ac+o(1)}$.
\end{theorem}

\begin{proof}
The prime number theorem gives $N\sim L^A/(A\log L)$, since the number
of primes dividing $q$ is $O(L)$. Equation~\eqref{eq:sharp} applies to every nonprincipal character. In particular,
\[
 \max_{\chi\ne1}b_\chi(P)\le \epsilon_L,
 \qquad \epsilon_L=L^{1-A/2+o(1)}.
\]

Set $j=\lfloor\log(q-1)/\log y\rfloor$. Each product of $j$ primes
from $P$, with repetitions allowed, is an integer strictly below $q$.
Orthogonality counts congruent pairs of such ordered products. Since
congruence here is equality, unique factorization bounds the number of
possible partners of each ordered tuple by $j!$. Therefore
\begin{equation}\label{eq:shortmoment}
 \sum_\chi b_\chi(P)^{2j}
 \le\frac{\varphi(q)j!}{N^j}
 \le q^{1/A+o(1)}.
\end{equation}
Indeed $j=(1/A+o(1))L/\log L$ and
$\log(j/N)=-(A-1)\log L+O(1)$.

Fix $c'$ strictly between $c_A$ and $c$, and let $s$ be the smallest even
integer at least $c'L/\log L$. Since $c_A>2/A$, we have $s>2j$.
Multiplying the nonprincipal terms of \eqref{eq:shortmoment} by
$\epsilon_L^{s-2j}$ bounds their $s$th moment by
\[
 q^{1/A-(A/2-1)(c'-2/A)+o(1)}=o(1).
\]
The strict inequality in the exponent follows from the definition of $c_A$.
For the even exponent $s$, character orthogonality counts ordered pairs of
$(s/2)$-tuples with congruent products, all with positive weight.
Every such pair in $S$ is a pair in $P$. Restoring the normalization gives
\begin{equation}\label{eq:momenttransfer}
 \sum_\chi b_\chi(S)^s
 \le (N/m)^s \sum_\chi b_\chi(P)^s
 \le a^{-s}(1+o(1)).
\end{equation}
The principal term $1$ is included before transferring this inequality.

Complementing subsets gives $|E_k(\chi)|=|E_h(\chi)|$.
If $h\ge BL$ for a sufficiently large fixed $B$, the circle bound from
the preceding sections gives
$q(m+1)\exp(-\tau h/2)$ for the desired sum, which is exponentially
small in $L$. Hence assume $cL/\log L\le h\le BL$.
For $u\ge b_\chi(S)$ and $r=h/(mu)>0$ one has
\[
 \prod_{p\in S}|1+z\chi(p)|
 \le\exp(mru+mr^2/2)\quad (|z|=r).
\]
This follows by applying $1+x\le e^x$ to each squared modulus,
then bounding the real part of $z\sum\chi(p)$ by $mru$.
Cauchy's inequality and $\binom mh\ge(m/h)^h$ give the convenient bound
\begin{equation}\label{eq:smallcoeff}
 |E_h(\chi)|\le u^h\exp\!\left(h+\frac{h^2}{2mu^2}\right)
 \le e^{3h/2}u^h
 \quad\text{if } u\ge u_0:=\sqrt{h/m}.
\end{equation}

Uniformly over this range of $h$,
$u_0=L^{-(A-1)/2+o(1)}$. For $b_\chi(S)\le u_0$, the
total contribution in \eqref{eq:smallcoeff} is at most
$q(e^{3/2}u_0)^h\le q^{1-c(A-1)/2+o(1)}$, exponentially small in
$L$ because $c_A>2/(A-1)$.
For $u_0<b_\chi(S)\le\zeta$, use \eqref{eq:smallcoeff} with
$u=b_\chi(S)$ and apply \eqref{eq:momenttransfer}. This gives at most
\[
 (1+o(1))e^{3h/2}\zeta^{h-s}a^{-s}.
\]
Here $h-s\ge\kappa h$ for some fixed $\kappa>0$ and all sufficiently
large $q$. Choosing $0<\zeta<1-\tau$ sufficiently small in terms of
$a,c,c'$ bounds this expression by $e^{-h}$.

By the Gram matrix bound, only $O_{a,\zeta}(1)$ characters have
$b_\chi(S)>\zeta$, since the off-diagonal inner products on $P$ are
at most $\epsilon_L=o(1)$. For these remaining characters outside $H$,
the circle bound gives $(m+1)\exp(-\tau h/2)$ each. Since $\log m=O(\log L)$,
their sum is exponentially small in $h$. Combining the three ranges proves
the estimate. The minimum positive probability in the quotient proved
above is an inverse power of $m$, so this error suffices for exact
support. Finally $y^k=\exp(Ak\log L)=q^{Ac+o(1)}$.
\end{proof}

For $a>1/2$, the quotient is trivial for all sufficiently large $q$.
Thus for each fixed $\eta>0$ every unit has a squarefree representative
at most $q^{2(A-1)/(A-2)+\eta}$ supported on any retained fraction
at least $a$. For example, $A=4$ gives the exponent $3+\eta$ using
only primes at most $(\log q)^4$. A counting argument determines the
correct order of the required number of factors: for prime $q$ and
$k=(b+o(1))L/\log L$ with fixed $b>0$,
\[
 \binom Nk=q^{b(A-1)+o(1)}.
\]
Indeed $\log\binom Nk=k\log(N/k)+O(k+k^2/N)$ by Stirling's formula.
Even the full prime pool has fewer than $q-1$ different $k$-subsets
when $b<1/(A-1)$. These bounds leave the optimal leading constant undetermined.
\section{An unconditional endpoint for almost all prime moduli}
\label{sec:large-sieve}

At a common cutoff for each interval of moduli the multiplicative large sieve
supplies the cancellation needed for the logarithm-squared result.  The
exceptional set can be estimated before the retained primes are selected.

\begin{theorem}\label{thm:density-endpoint}
Fix $0<a\le1$ and $b>0$. There are constants $B=B(a)>0$ and $C=C(a,b)>0$
with the following property. Put $y=C(\log Q_0)^2$. For all sufficiently
large $Q_0$, apart from at most
\[
 Q_0\exp\!\left(-b\frac{\log Q_0}{\log\log Q_0}\right)
\]
prime moduli $q\in[Q_0,2Q_0]$, the quotient bound, error estimate and exact
support of Theorem~\ref{thm:main} hold for every $S$ of density at least
$a$ and every $\min(k,m-k)\ge B\log q$. This assertion is unconditional.
\end{theorem}

\begin{proof}
Temporarily write $Q=Q_0$, $L=\log Q$, $\ell=\log L$, and $N=\pi(y)$.
For large $Q$ we have $y<Q$ and all these primes are eligible.
Fix a small $\xi>0$ depending only on $a$.
Call a prime $q\in[Q,2Q]$ exceptional if it has a nonprincipal character
$\chi$ for which $|\sum_{p\le y}\chi(p)|>\xi N$.
Put $j=\lfloor 2L/\log y\rfloor$, and let $a_n$ count ordered
$j$-tuples of primes up to $y$ with product $n$.
Unique factorization gives $0\le a_n\le j!$, and hence
$\sum_n a_n^2\le j!N^j$. Since $y^j\le Q^2$, the classical
multiplicative large sieve gives
\[
 \sum_{d\le2Q}\frac{d}{\varphi(d)}
 \sum_{\chi\bmod d}^{*}\left|\sum_{n\le y^j}a_n\chi(n)\right|^2
 \le (4Q^2+y^j)\sum_n a_n^2 .
\]
Here the star restricts the sum to primitive characters.  See
D.~Koukoulopoulos, \emph{The Distribution of Prime Numbers},
Theorem~25.15, for this inequality with an absolute implied constant,
which suffices throughout; the displayed normalization also follows from
the sharp additive large sieve by the primitive Gauss sum identity.
Every nonprincipal character modulo a prime is primitive.
Each exceptional prime therefore contributes at least $(\xi N)^{2j}$
to the left. If $E$ is their number, it follows that
\[
 E\le 5Q^2\frac{j!}{(\xi^2N)^j}.
\]

To retain the gain at the critical exponent, we keep the next term.
Since $N\sim y/\log y$ and $j=2L/(2\ell+\log C)+O(1)$, we have
$\log(N/j)=\ell+\log(C/2)+o(1)$.  Stirling's formula gives
\begin{align*}
 \log\!\left(5Q^2\frac{j!}{(\xi^2N)^j}\right)
 &= 2L-j\bigl(\ell+\log(C/2)+1+2\log\xi+o(1)\bigr)
      + O(\log j)\\
 &=L-\left(\tfrac12\log C+1+2\log\xi-\log 2+o(1)\right)
           \frac{L}{\ell}.
\end{align*}
Choose $C$ so that the coefficient in parentheses without $o(1)$
exceeds $b$. This proves the claimed exceptional-set bound.

For every remaining modulus, all full-pool nonprincipal biases are at most
$\xi$. Take $\xi$ at most the coherence tolerance chosen in the proof of
Theorem~\ref{thm:main}. The Gram bound, circle estimate and swap lemma
then give all the conclusions with $B=10/\tau$, where $\tau$ is the
character threshold used in that theorem. There are asymptotically
$Q/\log Q$ primes in the interval, so the exceptional fraction tends to zero.
\end{proof}

\begin{corollary}[Shorter products for almost every prime modulus]
Fix $A>2$, $0<a\le1$ and $c>2/(A-1)$. With $y=(\log Q_0)^A$, the
conclusions of Theorem~\ref{thm:short-products} hold unconditionally for
all but $O(Q_0^{1-\eta})$ prime moduli $q\in[Q_0,2Q_0]$, for some
$\eta>0$. Here $c$ replaces the constant in the length threshold.
\end{corollary}
\begin{proof}
Put $L=\log Q_0$ and choose $2/(A-1)<c'<\min(c,4/A)$, which is possible
since $A>2$. Take $s=2j$ even with $s=c'L/\log L+O(1)$. Then
$y^j\le Q_0^2$ for large $Q_0$. The same large sieve gives
\[
 \sum_{\substack{Q_0\le q\le2Q_0\\ q\ \mathrm{prime}}}
 \ \sum_{\chi\ne1} b_\chi(P)^s
 \ll Q_0^2 j! /N^j=Q_0^{2-c'(A-1)/2+o(1)} .
\]
The exponent is strictly below one. Choose $\eta>0$ small enough that
the right side is $O(Q_0^{1-2\eta})$. Except for $O(Q_0^{1-\eta})$ prime
moduli the inner sum is at most $Q_0^{-\eta}$. Thus the moment including
the principal character is $1+o(1)$, and each nonprincipal bias is at
most $Q_0^{-\eta/s}=o(1)$. These were the two full-pool estimates used
in the short-product proof. Apply \eqref{eq:momenttransfer} and
\eqref{eq:smallcoeff} with this $s$. The low-bias contribution tends to
zero exponentially in $L$ since $c>2/(A-1)$, the middle range uses
$c>c'$, and the Gram and circle bounds handle the remaining characters.
The substitution $\log q=L+O(1)$ does not affect these strict inequalities.
\end{proof}
\section{One logarithmic-size list of primes for many moduli}
\begin{corollary}\label{cor:common-list}
Fix $0<\delta<1/2$ and $b>0$. There are constants $C,D>0$ and, for
each sufficiently large $Q$, a single list $U$ of $M=\lceil D\log Q\rceil$
distinct primes at most $C(\log Q)^2$ with the following property.
Under Dirichlet GRH one list works simultaneously for all integers
$q\in[Q,2Q]$. There is also, unconditionally, a single list working
simultaneously for all primes in that interval outside a set of at most
$Q\exp(-b\log Q/\log\log Q)$ exceptions.
Write $U_q=\{p\in U:p\nmid q\}$ and $k=\lfloor M/4\rfloor$.
For every $S\subseteq U_q$ of size $m\ge(1/2+\delta)|U_q|$, the
number of $k$-subsets giving any specified unit modulo $q$ is
$\binom mk(1+O(Q^{-2}))/\varphi(q)$.
\end{corollary}

\begin{proof}
Put $y=C(\log Q)^2$ and $\gamma=\delta/8$.
For all relevant $q$, the mean of each nonprincipal $\chi$ over
\emph{all} primes up to $y$, using $\chi(p)=0$ on nonunits, has modulus
at most $\gamma$. Under GRH this follows from the proof of
Lemma~\ref{lem:analytic}, uniformly since $\log q\le\log(2Q)$.
Unconditionally use the exceptional-set estimate in the proof of
Theorem~\ref{thm:density-endpoint} with $\xi=\gamma$.
Choose $U$ uniformly among the $M$-subsets of these primes. Hoeffding's
inequality \cite{hoeffding} for sampling without replacement, applied to real and
imaginary parts in $[-1,1]$, bounds the probability that any one
sample mean differs from its population mean by more than $\gamma$ by
$4\exp(-M\gamma^2/4)$. There are at most $\sum_{q\le 2Q}\varphi(q)
=O(Q^2)$ characters, so a sufficiently large $D$ guarantees a list
with every required sample bias at most $2\gamma$.

For this list $|U_q|=M-o(M)$ uniformly, since $\omega(q)!\le q$
implies $\omega(q)=O(\log Q/\log\log Q)$. Thus
$b_\chi(U_q)\le3\gamma$ and deletion gives
\[
 b_\chi(S)\le\frac{1/2-\delta+3\gamma}{1/2+\delta}
 =:1-\kappa<1 .
\]
For large $Q$ we have $h=\min(k,m-k)\ge M/5$.
By \eqref{eq:circle}, summing over nonprincipal characters costs at most
$2Q(M+1)e^{-\kappa M/10}$. Choosing $D$ large enough makes this
$O(Q^{-2})$. Fourier inversion proves the assertion.
\end{proof}
Even with unrestricted subset sizes, covering all units modulo a prime
$q\in[Q,2Q]$ requires $2^{|U|}\ge q-1$. Thus the order $\log Q$
of the common list is optimal.
\section{Finite computations}
The accompanying Python program enumerates all retained prime sets of the
indicated sizes. For each it computes the counts by direct enumeration of
subsets and independently by a recurrence on the residue and number of primes
used. Descending updates in the cardinality coordinate enforce distinctness.
The last column is the minimum number of representations over every retained
set of that size and every unit.
\begin{center}
\begin{tabular}{rrrrrrr}
\toprule
$q$&$y$&$N$&$m$&$k$&sets missing a unit / all sets&minimum\\
\midrule
7&47&14&8&3&1 / 3003&0\\
7&47&14&9&3&0 / 2002&5\\
15&37&10&6&3&21 / 210&0\\
15&37&10&7&3&0 / 120&2\\
\bottomrule
\end{tabular}
\end{center}
For $q=15$ exactly five of the ten primes are $+1$ under the quadratic
character modulo $3$. Retaining those five gives ten positive-sign triples.
Adding two negative-sign primes gives fifteen positive and twenty negative
triples, as the quotient polynomial predicts. A cubic example shows a different phenomenon. Modulo $13$, retain the
thirteen primes up to $211$ in the cube subgroup and the single additional
prime $2$. Five-factor products occupy eight of the twelve units. Their
counts in the three cubic classes are $1287,715,0$, even though both
quadratic signs occur in the retained set. The numbers come from
$\binom{13}{5}$ and $\binom{13}{4}$.

\section{What the three starting sources establish}

The formal input is Dirichlet GRH for every positive modulus.  In the
readable edition of \texttt{lsernau/grh-lean}, the statement file imports
only Mathlib.  Its half-plane assertion quantifies over every Dirichlet
character, including imprimitive characters, and excludes the principal
character at its pole.  Its critical-strip assertion puts every zero with
$0<\Re s<1$ on $\Re s=1/2$.  The main file exports both assertions and
Mathlib's \texttt{RiemannHypothesis}.  This fixes exactly which family of
zeros supplies the character estimates in our proof.
\href{https://github.com/lsernau/grh-lean/blob/72fead102f663d5c7a407964a13dc0cf9d977207/grh/src/Statements.lean}{Statement file},
\href{https://github.com/lsernau/grh-lean/blob/72fead102f663d5c7a407964a13dc0cf9d977207/grh/src/Main.lean}{exported theorems}.

Alpöge's argument in this edition bounds the same corrected lattice
covolume from above and below.  Cancellation of their common main term
bounds a one-sided weighted prime-sum remainder, with kernel $1/n-1/x$
for $n\le x$.  The zero criterion turns this bound into half-plane
nonvanishing.  Here we use the zero theorem
to control averages of characters over small primes, then count products
after deletions.  The project pins Lean 4.35.0 and Mathlib
\texttt{56a00585}.  Its maintainers report that the five exported theorems
depend on \texttt{propext}, \texttt{Classical.choice}, and
\texttt{Quot.sound}, and report successful checks by three independent
kernel implementations.  Our verification of the upstream input consisted
of inspecting the pinned source declarations and build documentation.
\href{https://github.com/lsernau/grh-lean/blob/72fead102f663d5c7a407964a13dc0cf9d977207/grh/src/New/Reduction/MainInequality.lean}{Covolume comparison},
\href{https://github.com/lsernau/grh-lean/blob/72fead102f663d5c7a407964a13dc0cf9d977207/grh/src/New/Reduction/OneSidedBound.lean}{one-sided reduction},
\href{https://github.com/lsernau/grh-lean/blob/72fead102f663d5c7a407964a13dc0cf9d977207/grh/README.md}{build and verification record}.

The ProveIt directory indexes 61 research packets and describes their
placement as archival.  Two illustrate the difference in the question
being asked.  The cyclic-generator packet uses signed Frobenius weights
to test zeros of a Dedekind zeta function.  The sparse prime-sampling
packet uses smoothed sums to recover zero multiplicities.  Here the
observable is the residue of a product of a prescribed number of distinct
primes.  The exceptional characters describe the distribution of that
product on a small quotient.
\href{https://github.com/VladimirReshetnikov/ProveIt/blob/4d7990202f9f106b2ffe8c209f073ad01d5e1987/Analysis/GeneralizedRiemannHypothesis/README.md}{ProveIt inventory},
\href{https://github.com/VladimirReshetnikov/ProveIt/blob/4d7990202f9f106b2ffe8c209f073ad01d5e1987/Analysis/GeneralizedRiemannHypothesis/Research/grh_cyclic_generator_transfer_2026-11-25/README.md}{cyclic-generator packet},
\href{https://github.com/VladimirReshetnikov/ProveIt/blob/4d7990202f9f106b2ffe8c209f073ad01d5e1987/Analysis/GeneralizedRiemannHypothesis/Research/grh-torsion-and-prime-sampling/README.md}{prime-sampling packet}.

The supplied encyclopedia article records the release of Alpöge's
formalization on 20 November 2026 and explains the classical relationship
between the real parts of zeta zeros and the error in the prime number
theorem.  For the residue problem this relationship must hold after
twisting by each character modulo $q$.  Ordinary RH supplies the principal
case, and the Dirichlet GRH supplies the nonprincipal cases.
\href{https://en.wikipedia.org/wiki/Riemann_hypothesis}{Riemann hypothesis,
sections on the formal proof and the distribution of primes, consulted
26 November 2026}.

The closest earlier report available for this project is
\emph{Squarefree products after adversarial deletion of small primes}
(26 November 2026, in \texttt{orthogonal\_nullity\_and\_prime\_products.zip}).
It gives the majority result for $y=L^A$, $A>2$, at $h\gg\log q$, and
bounds individual characters according to their order. We reuse its
fixed-cardinality circle argument in \eqref{eq:circle}. The short-word
argument here bounds the order of the entire exceptional subgroup, the
swap lemma recovers exact support even on rare fibers, and the moment
transfer reduces the required number of factors. Smoothing gives the
separate $CL^2$ endpoint.
\section{Relation to the literature and further questions}
Booker--Pomerance \cite{bp} distinguished the size of a representative
from the size of its largest prime factor. Munsch--Shparlinski \cite{ms}
proved the polynomial-height assertion for $p$-smooth representatives and
obtained representatives of height $p^{2+o(1)}$ at smoothness
$p^{1/(4\sqrt e)+o(1)}$. Martin--Parvardi \cite{mp} studied subproducts of
all small integers. Jao--Miller--Venkatesan \cite{jmv} used GRH to obtain
expansion from small primes. Matom\"aki--Ter\"av\"ainen \cite{mt} studied
products of a bounded number of much larger primes and identified the role
of small-index cosets in that setting.

Klurman--Shparlinski--Ter\"av\"ainen \cite{kst} explicitly observed under GRH
that unrestricted subset products from all primes up to $(\log q)^B$, $B>4$,
cover every unit. Their main almost-all theorem constructs a common small
set of primes and controls the number of exceptional moduli. Hughes also studies
common prime lists of asymptotically optimal cardinality for almost all moduli
(\href{https://arxiv.org/abs/2610.04678}{arXiv:2610.04678}). The results
here give an endpoint cutoff, exact cardinalities, arbitrary retained sets,
and relative counts through a bounded quotient. The comparison identifies
closely related methods and results. Publication priority for the combined
statements has not been established.
\begin{enumerate}
\item At $y=C\log^2 q$, can every unit be reached after majority retention
with $O(\log q/\log\log q)$ prime factors? This would make the corresponding
representatives polynomially bounded in $q$ at the proposed smoothness scale.
\item For $A>2$, what is the optimal constant in front of
$\log q/\log\log q$ for uniform relative counts? The moment bound and the
elementary counting obstruction leave a gap.
\item Can the finite quotient be computed from the retained primes in time
polynomial in $\log q$ without enumerating $\widehat G$? Its bounded order
makes the subsequent quotient counting inexpensive.
\end{enumerate}

\begin{thebibliography}{10}
\bibitem{bp} A. R. Booker and C. Pomerance, \emph{Squarefree smooth numbers
and Euclidean prime generators}, Proc. Amer. Math. Soc. 145 (2017),
5035--5042, especially \S6. \href{https://arxiv.org/abs/1607.01557}{arXiv:1607.01557}.
\bibitem{mv} H. L. Montgomery and R. C. Vaughan, \emph{Multiplicative Number
Theory I. Classical Theory}, Cambridge, 2007. Theorem 13.7 and proof of
Theorem 13.11, equations (13.21) and (13.28).
\href{https://doi.org/10.1017/CBO9780511618314.015}{Chapter 13}.
\bibitem{ms} M. Munsch and I. E. Shparlinski, \emph{On smooth square-free
numbers in arithmetic progressions}, J. London Math. Soc. 101 (2020).
\href{https://arxiv.org/abs/1710.04705}{arXiv:1710.04705}.
\bibitem{mp} G. Martin and A. Parvardi, \emph{Subproducts of small residue
classes}, Canad. Math. Bull. 65 (2022), 1--8.
\href{https://arxiv.org/abs/2008.10198}{arXiv:2008.10198}.
\bibitem{jmv} D. Jao, S. D. Miller and R. Venkatesan,
\emph{Expander graphs based on GRH with an application to elliptic curve
cryptography}, J. Number Theory 129 (2009).
\href{https://arxiv.org/abs/0811.0647}{arXiv:0811.0647}.
\bibitem{mt} K. Matom\"aki and J. Ter\"av\"ainen, \emph{Products of primes in
arithmetic progressions}. \href{https://arxiv.org/abs/2301.07679}{arXiv:2301.07679}.
\bibitem{kst} O. Klurman, I. E. Shparlinski and J. Ter\"av\"ainen,
\emph{On the generation of multiplicative groups by small primes}, 2026.
\href{https://arxiv.org/abs/2609.25393}{arXiv:2609.25393}.
\bibitem{hoeffding} W. Hoeffding, \emph{Probability inequalities for sums
of bounded random variables}, J. Amer. Statist. Assoc. 58 (1963), 13--30,
especially \S6. \href{https://doi.org/10.1080/01621459.1963.10500830}{DOI}.
\end{thebibliography}
\end{document}
